% GENERATED FILE. Edit canonical lessons, then run npm run generate:books.
% canonical-source-hash: fnv1a64:e2a7ff156a709901
% canonical-lessons: GU-C186-cukavvu, GU-C186-rokad, GU-C186-parcuran, GU-C186-not, GU-C186-rasid

\chapter{Paying at the till}
\label{ch:gu-paying-at-the-till}
\hlchaptermodality{186}

\begin{chapteropening}
\textbf{By the end of this chapter:} \emph{I can pay in cash or in small change and get a receipt.}

Complete the last lesson of chapter 186: I can pay in cash or in small change and get a receipt.
\end{chapteropening}

\section[\texorpdfstring{\gu{ચૂકવવું} (cūkavvũ)}{ચૂકવવું (cūkavvũ)}]{\gu{ચૂકવવું} (cūkavvũ) — to pay, to settle}

\label{lesson:GU-C186-cukavvu}



\begin{quote}
Before the new one: say the Gujarati for pair, then the Gujarati for to test.
\end{quote}

\subsection*{You'll want to know: \gu{ચૂકવવું}}
\textbf{\gu{ચૂકવવું}} — \emph{cūkavvũ} — \textquotedblleft{}to pay, to settle\textquotedblright{}.

\begin{quote}
\textbf{\gu{હું} \gu{દુકાનદારને} \gu{પૈસા} \gu{ચૂકવું} \gu{છું}.} — \emph{hũ dukāndārne paisā cūkavũ chũ} — \textquotedblleft{}I pay the shopkeeper the money.\textquotedblright{}
\end{quote}

\begin{cousinweb}
\textbf{\gu{ચૂકવવું}} (\emph{cūkavvũ}) is a tadbhava, the causative of a verb built on \emph{*cukk}, \textquotedblleft{}fall short of, stop\textquotedblright{}, a reconstructed form. CDIAL glosses the Gujarati word \textquotedblleft{}to discharge, settle\textquotedblright{} (CDIAL 4848).
\end{cousinweb}

\begin{grammarlens}[title={What you've built}]
One more word for everyday life at A2.
\end{grammarlens}

\subsection*{Your turn}
\begin{itemize}
\raggedright
  \item \emph{Say it:} \emph{cūkavvũ}
  \item \emph{Say it:} \emph{cūkavvũ}, once more
  \item \emph{Say it:} \emph{hũ dukāndārne paisā cūkavũ chũ} — I pay the shopkeeper the money.
  \item \emph{Recall:} say the Gujarati for pair, then the Gujarati for to test, then say \emph{cūkavvũ} again
\end{itemize}

\begin{tcolorbox}[breakable,colback=teal!4,colframe=teal!35!black,title={Before you move on}]
What does \gu{ચૂકવવું} mean? (\textquotedblleft{}To pay, to settle\textquotedblright{}.) Say it once more.
\end{tcolorbox}
\section[\texorpdfstring{\gu{રોકડ} (rokaḍ)}{રોકડ (rokaḍ)}]{\gu{રોકડ} (rokaḍ) — cash}

\label{lesson:GU-C186-rokad}



\begin{quote}
Before the new one: say the Gujarati for to test, then the Gujarati for to pay.
\end{quote}

\subsection*{You'll want to know: \gu{રોકડ}}
\textbf{\gu{રોકડ}} — \emph{rokaḍ} — \textquotedblleft{}cash\textquotedblright{}.

\begin{quote}
\textbf{\gu{મારી} \gu{પાસે} \gu{રોકડ} \gu{નથી}.} — \emph{mārī pāse rokaḍ nathī} — \textquotedblleft{}I have no cash on me.\textquotedblright{}
\end{quote}

A feminine noun.

\begin{cousinweb}
\textbf{\gu{રોકડ}} (\emph{rokaḍ}) is a tadbhava. CDIAL lists it, glossed \textquotedblleft{}cash\textquotedblright{}, under \emph{*rōkka}, \textquotedblleft{}ready money\textquotedblright{}, a reconstructed form, extended with the suffix \emph{-ḍ-} (CDIAL 10828).
\end{cousinweb}

\begin{grammarlens}[title={What you've built}]
One more word for everyday life at A2.
\end{grammarlens}

\subsection*{Your turn}
\begin{itemize}
\raggedright
  \item \emph{Say it:} \emph{rokaḍ}
  \item \emph{Say it:} \emph{rokaḍ}, once more
  \item \emph{Say it:} \emph{mārī pāse rokaḍ nathī} — I have no cash on me.
  \item \emph{Recall:} say the Gujarati for to test, then the Gujarati for to pay, then say \emph{rokaḍ} again
\end{itemize}

\begin{tcolorbox}[breakable,colback=teal!4,colframe=teal!35!black,title={Before you move on}]
What does \gu{રોકડ} mean? (\textquotedblleft{}Cash\textquotedblright{}.) Say it once more.
\end{tcolorbox}
\section[\texorpdfstring{\gu{પરચૂરણ} (parcūraṇ)}{પરચૂરણ (parcūraṇ)}]{\gu{પરચૂરણ} (parcūraṇ) — small change; odds and ends}

\label{lesson:GU-C186-parcuran}



\begin{quote}
Before the new one: say the Gujarati for to pay, then the Gujarati for cash.
\end{quote}

\subsection*{You'll want to know: \gu{પરચૂરણ}}
\textbf{\gu{પરચૂરણ}} — \emph{parcūraṇ} — \textquotedblleft{}small change; odds and ends\textquotedblright{}.

\begin{quote}
\textbf{\gu{દુકાનદાર} \gu{પાસે} \gu{પરચૂરણ} \gu{નથી}.} — \emph{dukāndār pāse parcūraṇ nathī} — \textquotedblleft{}The shopkeeper has no small change.\textquotedblright{}
\end{quote}

A neuter noun.

\begin{cousinweb}
The origin of \textbf{\gu{પરચૂરણ}} (\emph{parcūraṇ}) is not settled. CDIAL has no entry for it, and EtymDB lists no source.
\end{cousinweb}

\begin{grammarlens}[title={What you've built}]
One more word for everyday life at A2.
\end{grammarlens}

\subsection*{Your turn}
\begin{itemize}
\raggedright
  \item \emph{Say it:} \emph{parcūraṇ}
  \item \emph{Say it:} \emph{parcūraṇ}, once more
  \item \emph{Say it:} \emph{dukāndār pāse parcūraṇ nathī} — The shopkeeper has no small change.
  \item \emph{Recall:} say the Gujarati for to pay, then the Gujarati for cash, then say \emph{parcūraṇ} again
\end{itemize}

\begin{tcolorbox}[breakable,colback=teal!4,colframe=teal!35!black,title={Before you move on}]
What does \gu{પરચૂરણ} mean? (\textquotedblleft{}Small change; odds and ends\textquotedblright{}.) Say it once more.
\end{tcolorbox}
\section[\texorpdfstring{\gu{નોટ} (noṭ)}{નોટ (noṭ)}]{\gu{નોટ} (noṭ) — banknote}

\label{lesson:GU-C186-not}



\begin{quote}
Before the new one: say the Gujarati for cash, then the Gujarati for small change.
\end{quote}

\subsection*{You'll want to know: \gu{નોટ}}
\textbf{\gu{નોટ}} — \emph{noṭ} — \textquotedblleft{}banknote\textquotedblright{}.

\begin{quote}
\textbf{\gu{આ} \gu{સો} \gu{રૂપિયાની} \gu{નોટ} \gu{છે}.} — \emph{ā so rūpiyānī noṭ che} — \textquotedblleft{}This is a hundred-rupee note.\textquotedblright{}
\end{quote}

\begin{cousinweb}
\textbf{\gu{નોટ}} (\emph{noṭ}) is a loan from English \emph{note}, short for \emph{banknote}.
\end{cousinweb}

\begin{grammarlens}[title={What you've built}]
One more word for everyday life at A2.
\end{grammarlens}

\subsection*{Your turn}
\begin{itemize}
\raggedright
  \item \emph{Say it:} \emph{noṭ}
  \item \emph{Say it:} \emph{noṭ}, once more
  \item \emph{Say it:} \emph{ā so rūpiyānī noṭ che} — This is a hundred-rupee note.
  \item \emph{Recall:} say the Gujarati for cash, then the Gujarati for small change, then say \emph{noṭ} again
\end{itemize}

\begin{tcolorbox}[breakable,colback=teal!4,colframe=teal!35!black,title={Before you move on}]
What does \gu{નોટ} mean? (\textquotedblleft{}Banknote\textquotedblright{}.) Say it once more.
\end{tcolorbox}
\section[\texorpdfstring{\gu{રસીદ} (rasīd)}{રસીદ (rasīd)}]{\gu{રસીદ} (rasīd) — receipt}

\label{lesson:GU-C186-rasid}



\begin{quote}
Before the new one: say the Gujarati for small change, then the Gujarati for banknote.
\end{quote}

\subsection*{You'll want to know: \gu{રસીદ}}
\textbf{\gu{રસીદ}} — \emph{rasīd} — \textquotedblleft{}receipt\textquotedblright{}.

\begin{quote}
\textbf{\gu{દુકાનદાર} \gu{રસીદ} \gu{આપે} \gu{છે}.} — \emph{dukāndār rasīd āpe che} — \textquotedblleft{}The shopkeeper gives a receipt.\textquotedblright{}
\end{quote}

A feminine noun.

\begin{cousinweb}
\textbf{\gu{રસીદ}} (\emph{rasīd}) is a loan, usually traced to Persian \emph{rasīd}, \textquotedblleft{}receipt\textquotedblright{}; EtymDB gives the same Persian source. CDIAL has no entry for it.
\end{cousinweb}

\begin{grammarlens}[title={What you've built}]
One more word for everyday life at A2.
\end{grammarlens}

\subsection*{Your turn}
\begin{itemize}
\raggedright
  \item \emph{Say it:} \emph{rasīd}
  \item \emph{Say it:} \emph{rasīd}, once more
  \item \emph{Say it:} \emph{dukāndār rasīd āpe che} — The shopkeeper gives a receipt.
  \item \emph{Recall:} say the Gujarati for small change, then the Gujarati for banknote, then say \emph{rasīd} again
\end{itemize}

\begin{tcolorbox}[breakable,colback=teal!4,colframe=teal!35!black,title={Before you move on}]
What does \gu{રસીદ} mean? (\textquotedblleft{}Receipt\textquotedblright{}.) Say it once more.
\end{tcolorbox}
